%%%%%%%%%%  scan idx 188  |  folio 175  %%%%%%%%%%
\SGAPrintPage{175}%
\uline{are fully faithful and have as essential image} $\mathrm D^b(S)_{\coh}$ (I 2.15)\uline{. Moreover, when} $S$ \uline{has \textquotedblleft finite cohomological dimension\textquotedblright{} in the sense of} (3.7)\SGAAuthorFootnote{1}{Section~3 is independent of the rest of the exposé.} \uline{(for example, when the underlying space of} $S$ \uline{is noetherian and of finite Krull dimension), the functors}
\[
\mathrm K^{-}(\underline{L}(S))/\mathrm K^{-,\varnothing}(\underline{L}(S))\longrightarrow
\mathrm K^{-}(\Loclib(S))/\mathrm K^{-,\varnothing}(\Loclib(S))\longrightarrow
\mathrm D^{-}(S)_{\coh}
\]
\uline{are equivalences of categories.}
\item \uline{The canonical functor}
\[
\mathrm K^b(\Loclib(S))/\mathrm K^{b,\varnothing}(\Loclib(S))\longrightarrow\mathrm D(S)
\]
\uline{is fully faithful and has as essential image} $\mathrm D(S)_{\parf}$ (I 4.9)\uline{. For} $E\in\Ob\mathrm D(S)$ \uline{to have perfect amplitude contained in} $[m,n]$ (I 4.8)\uline{, it is necessary and sufficient that} $E$ \uline{be isomorphic, in} $\mathrm D(S)$\uline{, to a complex} $F$ \uline{such that} $F^i=0$ \uline{for} $i\notin[m,n]$, $F^i\in\Ob \underline{L}(S)$ \uline{for} $i>m$, \uline{with} $F^m$ \uline{locally free of finite type.}
\end{enumerate}

\uline{Proof.} In view of (3.5) and (3.7)\SGAAuthorFootMark{1}, this is an immediate corollary of (2.2.8): the only point to verify, which follows trivially from (3.5) (resp. (3.7)), is that, for an object $E$ of $\mathrm D^b(\Qcoh(S))$ (resp. $\mathrm D(\Qcoh(S))$ when $S$ has finite cohomological dimension in the sense of (3.7)), the notions of pseudo-coherence and perfection of (2.2.1) coincide with the usual notions (I 2.15 and 4.9) when $E$ is regarded as an object of $\mathrm D(S)$.

\sgaStableTarget{sga6.e02.rem.2-2-10}\textbf{Remark 2.2.10.} a) When $S$ is an affine scheme with ring $A$, the family consisting of the sheaf $\underline{O}_S$ is ample. For this family, we have $\underline{L}(S)=\Lib(S)$. It is also clear that the categories $\mathrm K^{-,\varnothing}(\Lib(S))$ and $\mathrm K^{-,\varnothing}(\Loclib(S))$ are reduced to zero objects. We therefore conclude from (2.2.9) that, writing $\Mod(A)$ for the category of $A$-modules and $\mathrm D(A)=\mathrm D(\Mod(A))$, the functor $\underline{\underline{R}}\Gamma_S:\mathrm D^+(S)\longrightarrow\mathrm D^+(A)$ induces a fully faithful functor $\mathrm D^b(S)_{\coh}\longrightarrow\mathrm D(A)$ (resp. $\mathrm D(S)_{\parf}\longrightarrow\mathrm D(A)$) whose essential image is
%%%%%%%%%%  scan idx 189  |  folio 176  %%%%%%%%%%
\SGAPrintPage{176}%
the
triangulated subcategory $\mathrm D^b(A)_{\coh}$ (resp. $\mathrm D(A)_{\parf}$) formed by complexes with bounded cohomology and pseudo-coherent (resp. perfect) relative to the subcategory $\Lib(A)$ (resp. $\Loclib(A)$) of free (resp. projective) modules of finite type. We have commutative diagrams of equivalences of categories
\[
\begin{tikzcd}
\mathrm K^{-,b}(\Lib(S)) \arrow[r] \arrow[d,"\Gamma_S"] &
\mathrm K^{-,b}(\Loclib(S)) \arrow[r] \arrow[d,"\Gamma_S"] &
\mathrm D^b(S)_{\coh} \arrow[d,"\mathbb R\Gamma_S"] \\
\mathrm K^{-,b}(\Lib(A)) \arrow[r] &
\mathrm K^{-,b}(\Loclib(A)) \arrow[r] &
\mathrm D^b(A)_{\coh},
\end{tikzcd}
\]
\[
\begin{tikzcd}
\mathrm K^b(\Loclib(S)) \arrow[r] \arrow[d,"\Gamma_S"] &
\mathrm D(S)_{\parf} \arrow[d,"\mathbb R\Gamma_S"] \\
\mathrm K^b(\Loclib(A)) \arrow[r] &
\mathrm D(A)_{\parf},
\end{tikzcd}
\]
(where, in the upper diagram, one may replace $b$ by $-$ when $S$ has finite cohomological dimension in the sense of (3.7)).

b) Let $S$ be a scheme. It follows from a) that a complex of $\underline{O}_S$-Modules which is pseudo-coherent and has bounded cohomology is locally isomorphic (in $\mathrm D(S)$) to a strictly pseudo-coherent complex, which was by no means evident a priori (cf. (I 2.3.1)).

\sgaStableTarget{sga6.e02.sec.2.3}\subsection*{2.3. Soft sheaves}

\sgaStableTarget{sga6.e02.para.2-3-1}2.3.1. Recall some definitions and results from [T.F]. Let $X$ be a paracompact space. A sheaf of sets $F$ on $X$ is said to be \uline{soft} when, for every closed subset $Y$ of $X$, every section of $F$ over $Y$ extends to a section of $F$ over $X$. Soft abelian sheaves are acyclic, i.e.\ when $F$ is soft, we have $H^i(X,F)=0$ for $i>0$ (T.F. II 4.4.3).

Let $A$ be a sheaf of rings on $X$. When $F$ is a soft $A$-Module, $F$ is generated by its sections (cf. 2.1.1): indeed, the homomorphism $A_X^{(I)}\to F$ defined by $I=H^0(X,F)$ is clearly an epimorphism.
%%%%%%%%%%  scan idx 190  |  folio 177  %%%%%%%%%%
\SGAPrintPage{177}%
On the other hand, when $A$ is soft, every $A$-Module is soft (T.F II 3.7.1). As examples of soft sheaves of rings, we mention the sheaf of continuous functions (real or complex) on a paracompact space, and the sheaf of $C^r$ functions on a manifold of class $C^r$.

\sgaStableTarget{sga6.e02.prop.2-3-2}\textbf{Proposition 2.3.2.} \uline{Let} $S$ \uline{be a compact space, equipped with a soft sheaf of rings} $A$\uline{. We write} $\Mod(S)$ \uline{for the category of left} $A$\uline{-Modules,} $\mathrm D(S)=\mathrm D(\Mod(S))$.
\begin{enumerate}[label=\alph*)]
\item \uline{The category} $\Lib(S)$ \uline{(resp.} $\underline{\operatorname{Loclib}}(S)$) (2.0) \uline{is liftable in} $\Mod(S)$ \uline{and quasi-stable} (resp. stable) \uline{under kernels of epimorphisms} (I 1.2)\SGAAuthorFootnote{1}{This assertion holds more generally for $S$ paracompact.}.
\item \uline{Let} $E\in\Ob\mathrm D(S)$ \uline{and} $n\in\mathbb Z$\uline{. For} $E$ \uline{to be} $n$\uline{-pseudo-coherent (resp.\ pseudo-coherent), it is necessary and sufficient that} $E$ \uline{be so relative to} $\Lib(S)$ \uline{or to} $\Loclib(S)$\uline{. The canonical functors}
\[
\mathrm K^{-}(\Lib(S))\longrightarrow \mathrm D(S)
\]
\uline{and}
\[
\mathrm K^{-}(\Loclib(S))\longrightarrow \mathrm D(S)
\]
\uline{are fully faithful and have as essential image} $\mathrm D^{-}(S)_{\coh}$ (I 2.15).
\item \uline{Let} $E\in\Ob\mathrm D(S)$\uline{, and} $[m,n]$ \uline{be an interval in} $\mathbb Z$\uline{. For} $E$ \uline{to have perfect amplitude contained in} $[m,n]$ (I 4.8), \uline{it is necessary and sufficient that there exist a quasi-isomorphism} $F\to E$\uline{, where} $F$ \uline{is a complex such that} $F^i=0$ \uline{for} $i\notin[m,n]$, $F^i\in\Ob\Lib(S)$ \uline{for} $i>m$ \uline{and} $F^m\in\Ob\Loclib(S)$\uline{. The canonical functor}
\[
\mathrm K^b(\Loclib(S))\longrightarrow\mathrm D(S)
\]
\uline{is fully faithful and has as essential image} $\mathrm D(S)_{\parf}$ (I 4.9).
\item \uline{Denote by} $\Mod(A(S))$ \uline{the category of left} $A(S)$\uline{-modules, and by} $\mathrm D(A(S))$ \uline{the corresponding derived category. The functor}
\[
H^0(S,.):\Mod(S)\longrightarrow\Mod(A(S))
\]
\uline{is exact and an equivalence of categories. Moreover, it
defines}
%%%%%%%%%%  scan idx 191  |  folio 178  %%%%%%%%%%
\SGAPrintPage{178}%
\uline{equivalences of categories:}
\[
\begin{gathered}
\mathrm D(S)\longrightarrow\mathrm D(A(S))\\
\mathrm D^{-}(S)_{\coh}\longrightarrow\mathrm D^{-}(A(S))_{\coh}\\
\mathrm D(S)_{\parf}\longrightarrow\mathrm D(A(S))_{\parf},
\end{gathered}
\]
\uline{the notions of pseudo-coherence (resp.\ perfection) in} $\mathrm D(A(S))$ \uline{being understood relative to the category of free (resp.\ projective) modules of finite type. It also defines an equivalence from the full subcategory of} $\mathrm D(S)$ \uline{formed by} $n$\uline{-pseudo-coherent objects (resp.\ objects of perfect amplitude contained in} $[m,n]$\uline{) onto the analogous full subcategory of} $\mathrm D(A(S))$.
\end{enumerate}

\uline{Proof.} The first assertion of a) follows from (2.1.2 a)) applied by taking for $\underline{S}$ the topos of sheaves of sets on $S$, and $\underline{C}'_S=\Mod(S)$ (in view of (2.3.1)). The second follows from (1.1 a)) (resp. (1.2 a))) with $\underline{C}=\Mod$ (the fibered category of $A$-Modules on variable open subsets), $\underline{C}_0=\Lib$ (resp. $\Loclib$), condition (ii) of (1.1) being verified by (2.3.1) and (2.1.1.b)). For assertion b), we apply (1.1 b),c)) (with $\underline{C}$ and $\underline{C}_0$ as before). Assertion c) follows from (1.2 b),c),d)), where $\underline{C}=\Mod$ and $\underline{C}_0=\Lib$. The first assertion of d) follows, by (2.3.1), from (2.1.2 b) (ii)), where we take $\underline{C}'_S=\underline{C}_S=\Mod(S)$. The other assertions of d) follow from the preceding one, in view of b) and c), and the fact that the functor $\Gamma_S$ sends free Modules (resp. direct factors of free Modules) of finite type to free modules (resp. projective modules of finite type).

\sgaStableTarget{sga6.e02.sec.2.4}\subsection*{2.4. Stein spaces}

\sgaStableTarget{sga6.e02.def.2-4-1}\textbf{Definition 2.4.1.} \uline{Let} $\underline{S}$ \uline{be a locally ringed topos} (I 2.15.1 (ii))\uline{, with final object} $S$\uline{. We say that} $\underline{S}$ \uline{is} Stein \uline{when the structure sheaf} $\underline{O}_S$ \uline{is coherent} (I 3.1) \uline{and the following properties hold:}

\textquotedblleft\uline{theorem} A\textquotedblright{} \uline{: every coherent} $\underline{O}_S$\uline{-Module} (I 3.1) \uline{is generated by its sections,}

\textquotedblleft\uline{theorem} B\textquotedblright{} \uline{: for every coherent} $\underline{O}_S$\uline{-Module} $F$\uline{, we have} $H^i(S,F)=0$ \uline{for} $i>0$.

%%%%%%%%%%  scan idx 192  |  folio 179  %%%%%%%%%%
\SGAPrintPage{179}%
\sgaStableTarget{sga6.e02.rem.2-4-1-1}\textbf{Remark 2.4.1.1.} When $\underline{S}$ has \textquotedblleft enough closed points\textquotedblright\SGAAuthorFootnote{1}{i.e.\ enough fiber functors (SGA~4 IV 6.4.1) of the type $i_s:s\to\underline{S}$, where $i_s$ is the inclusion of a closed subtopos of $\underline{S}$.} $s$ such that the Ideal $\underline{m}_s$ defined by $s$\SGAAuthorFootnote{2}{i.e.\ formed by sections of $\underline{O}_S$ whose fiber at $s$ belongs to the maximal ideal $m_s$ of $\underline{O}_{S,s}$.} is coherent, the relation \textquotedblleft$H^1(S,F)=0$ for every coherent sheaf $F$\textquotedblright{} implies theorem A. Indeed, let $F$ be a coherent $\underline{O}_S$-Module and $s$ a closed point of $S$. We have the exact sequence
\[
0\longrightarrow\underline{m}_sF\longrightarrow F\longrightarrow F/\underline{m}_sF\longrightarrow0,
\]
where $F/\underline{m}_sF$ is a coherent sheaf concentrated at $s$. Applying the functor $H^0(S,.)$ and writing the cohomology exact sequence, we find that $H^0(S,F)\longrightarrow H^0(S,F/\underline{m}_sF)$ is an epimorphism. Since $H^0(S,F/\underline{m}_sF)\simeq F_s/m_sF_s$, where $m_s$ is the maximal ideal of $\underline{O}_{S,s}$, we conclude by Nakayama.

The editor does not know whether the preceding implication is true in general (when one assumes only that $\underline{O}_S$ is a coherent sheaf of local rings).

The editor also does not know whether the relation \textquotedblleft$H^1(S,F)=0$ for every coherent sheaf $F$\textquotedblright{} implies theorem B.

\sgaStableTarget{sga6.e02.ex.2-4-2}\textbf{Examples 2.4.2.} A complex analytic Stein space [3], a compact convex subset of $\mathbb C^n$ equipped with the sheaf induced by the sheaf of holomorphic functions [3], the spectrum of a noetherian ring, and an affine rigid analytic space [13] define Stein topoi.

\sgaStableTarget{sga6.e02.lem.2-4-3}\textbf{Lemma 2.4.3.} \uline{Let} $\underline{S}$ \uline{be a Stein topos} (2.4.1)\uline{. We assume that the final object} $S$ \uline{is quasi-compact} (2.1.1)\uline{, and we denote by} $\Coh(S)$ \uline{the abelian category of coherent} $\underline{O}_S$\uline{-Modules, i.e.\ of finite presentation. Then:}
\begin{enumerate}[label=\alph*)]
\item \uline{The category} $\Lib(S)$ \uline{(resp.} $\Loclib(S)$) (2.0) \uline{is liftable in} $\Coh(S)$ \uline{and quasi-stable (resp.\ stable) under kernels of epimorphisms} (I 1.2).
\item \uline{The canonical functors}
\[
\mathrm K^{-}(\Lib(S))\longrightarrow\mathrm K^{-}(\Loclib(S))\longrightarrow\mathrm D^{-}(\Coh(S))
\]
\uline{are equivalences of categories.}
